\section{Theme Architecture}
\label{sec:themes}

A theme is a \emph{pure styling layer}: it accepts the geometry produced by the rendering
model (Section~\ref{sec:rendering}) and assigns colors, typography, shapes, and decorations,
producing a final visual specification. A theme operates exclusively on signals the renderer
has already computed---coordinates, status values, categories, progress fractions. It may not
alter geometry, and it may not require IR fields beyond those defined in
Section~\ref{sec:ir}~\cite{thinkcell,vegalite2017}.

\begin{quote}
\textbf{Theme-engine contract (binding):} A theme MUST accept all valid IR documents. A theme
MUST NOT require IR fields beyond the specification. Two valid themes applied to the same IR
may produce visually different outputs, but their underlying geometries---coordinates, stacking
order, label positions---must be identical.
\end{quote}

This section defines the complete theme schema, describes the five built-in themes, and
illustrates how the same IR renders visually differently under each.

\subsection{Theme Design Philosophy}
\label{sec:themes-philosophy}

Themes are built on four principles:

\begin{itemize}
  \item \textbf{Data-independent.} A theme knows nothing about the content of the IR. It maps
        semantic signals (status, category) to visual treatments through fixed lookup tables.
        No IR field other than those defined in Section~\ref{sec:ir} is read.
  \item \textbf{Composable.} Themes support single inheritance. A child theme specifies only
        the properties that differ from its parent; all others inherit. This enables branded
        variants (\texttt{acme-consulting} extending \texttt{consulting}) with minimal
        duplication.
  \item \textbf{Separable.} Swapping themes requires zero IR changes. The same YAML document
        renders correctly under any theme. Content and style are fully orthogonal.
  \item \textbf{Deterministic.} Themes introduce no layout variability. Given the same IR,
        theme, and rendering algorithm, output is fully determined (see
        Section~\ref{sec:rendering-determinism}).
\end{itemize}

\subsection{Theme Schema}
\label{sec:themes-schema}

The theme schema defines all visual properties a theme may configure. Properties not set by a
theme are inherited from its \texttt{extends} parent, defaulting to the built-in
\texttt{default} theme if no parent is specified. The schema is organized into blocks; all
property names, types, and defaults apply uniformly across themes.

\subsubsection{Canvas and Margins}
\label{sec:schema-canvas}

\begin{lstlisting}[language=yaml,caption={Theme schema: canvas block}]
canvas:
  width:            1200  # logical pixels; height is always computed
  background_color: "#FFFFFF"
  margin:
    top:    48            # px above the axis header
    right:  40            # px right of the last tick
    bottom: 48            # px below the last track row
    left:    0            # px left of the track header column
\end{lstlisting}

The canvas height is computed from track heights and inter-track gaps (\S\ref{sec:rendering-coords})
and is never specified in a theme.

\subsubsection{Typography}
\label{sec:schema-typography}

\begin{lstlisting}[language=yaml,caption={Theme schema: typography block}]
typography:
  font_family:         "Helvetica Neue, Arial, sans-serif"
  # Embedded font files required for deterministic metric computation:
  font_files:
    regular: "fonts/HelveticaNeue-Regular.woff2"
    bold:    "fonts/HelveticaNeue-Bold.woff2"
  font_size_base:    11  # pt; body text and activity labels
  font_size_axis:    10  # pt; axis tick labels
  font_size_title:   18  # pt; document title in canvas header
  font_size_subtitle: 13 # pt
  font_size_track:   11  # pt; track header labels
  font_weight_label:  600
  font_weight_axis:   400
  font_weight_header: 700
  locale_aware:      true # use metadata.locale for date label formatting
\end{lstlisting}

\paragraph{Font embedding requirement.}
Themes must specify embedded font files (WOFF2 or equivalent) for all families used in
layout-critical computations (label width, line height, character advance widths). System
fonts are permitted only for decorative text---e.g., HTML tooltip overlays---where exact
cross-platform metric agreement is not required. The renderer ships the embedded metrics
tables at build time.

\subsubsection{Axis Styling}
\label{sec:schema-axis}

\begin{lstlisting}[language=yaml,caption={Theme schema: axis block}]
axis:
  height:                40  # px; height of the time-axis header band
  tick_height:            6  # px; length of tick marks
  tick_label_offset:      4  # px; gap between tick mark base and label top
  gridline_color:    "#E0E0E0"
  gridline_width:       0.5  # px
  gridline_opacity:     0.8
  gridline_style:    "solid" # solid | dashed | none
  section_band_alpha:   0.04 # alternating column shading opacity (0 = off)
  today_marker:
    enabled:           false
    color:        "#CC2222"
    width:             1.5  # px
    style:           "solid" # solid | dashed
    label:           "Today"
    label_visible:    true
\end{lstlisting}

\subsubsection{Track Styling}
\label{sec:schema-track}

\begin{lstlisting}[language=yaml,caption={Theme schema: track block}]
track:
  header_width:         160  # px; fixed left column for track label text
  row_height:            48  # px; single-lane track height (no sub-lanes)
  sub_lane_height:       24  # px; height added per additional overlap lane
  max_sub_lanes:          8  # hard cap; excess activities go to last lane
  row_gap:                8  # px; vertical gap between consecutive tracks
  header_font_size:      11  # pt
  header_font_weight:   600
  header_color:    "#111111"
  header_background:   null  # null = transparent
  header_align:    "right"   # right | left | center
  separator_color: "#CCCCCC"
  separator_width:       0.5
  group_bracket_width:    4  # px; left-edge bracket for grouped tracks
  group_label_size:      10  # pt; group label above bracket
\end{lstlisting}

\subsubsection{Activity Styling}
\label{sec:schema-activity}

\begin{lstlisting}[language=yaml,caption={Theme schema: activity block}]
activity:
  bar_height:            24  # px; height within a single lane
  bar_radius:             2  # px; corner radius
  bar_top_margin:        12  # px; from lane top to bar top (centres bar in lane)
  min_width:              4  # px; minimum rendered bar width (see edge cases)
  tbd_extension_px:    null  # null = 2 * mean tick spacing in pixels
  approx_fade_px:         8  # px; gradient fade width for approximate-date edges
  label_inside_min_width: 40 # px; bar must be >= this to place label inside
  label_font_size:       10  # pt
  label_color_inside: "#FFFFFF"
  label_color_outside: "#111111"
  label_truncate_chars:  32  # max chars before ellipsis truncation
  progress_bar_height:    4  # px; filled strip at bar bottom for progress value
  progress_bar_alpha:    0.35
  progress_label_visible: false
\end{lstlisting}

\subsubsection{Milestone Styling}
\label{sec:schema-milestone}

\begin{lstlisting}[language=yaml,caption={Theme schema: milestone block}]
milestone:
  diamond_size:          14  # px; half-diagonal of the diamond polygon
  diamond_stroke_width: 1.5  # px
  cross_track_label_scale: 1.2  # label font size multiplier for cross-track
  label_offset_x:         8  # px; gap from diamond right vertex to label left
  label_font_size:        9  # pt
  label_max_width:       80  # px; truncate label if wider
  label_stack_offset:    14  # px; y-shift per collision pass (Phase 6)
  stack_offset_y:        20  # px; y-offset per stacking slot (Phase 4)
\end{lstlisting}

\subsubsection{Annotation and Legend Styling}
\label{sec:schema-annotation}

\begin{lstlisting}[language=yaml,caption={Theme schema: annotation and legend blocks}]
annotation:
  font_size:             9   # pt
  background_color: "#FFFFF0"
  background_alpha:      0.9
  border_color:    "#CCCC88"
  border_width:          0.5
  padding:               6   # px; all sides
  connector_style: "elbow"   # elbow | straight | none

legend:
  position:         "bottom" # bottom | right | top | none
  item_height:          14   # px
  item_gap:              8   # px
  font_size:             9   # pt
  swatch_size:          10   # px square
  column_spacing:       16   # px between columns
\end{lstlisting}

\subsubsection{Status-to-Style Map}
\label{sec:status-style-map}

Every theme must define a complete \texttt{status\_map} covering all seven IR status values.
The map is a lookup table: \texttt{status} $\to$ \{\,fill, stroke, opacity, pattern\,\}.

\begin{table}[h]
\centering
\small
\begin{tabular}{llllp{3.8cm}}
\toprule
\textbf{Status} & \textbf{Fill role} & \textbf{Opacity} & \textbf{Pattern} & \textbf{Visual intent} \\
\midrule
\texttt{planned}     & Accent-light  & 1.0  & \texttt{solid}          & Default; not yet started \\
\texttt{in-progress} & Accent-dark   & 1.0  & \texttt{solid}          & Active; full saturation \\
\texttt{done}        & Grey-mid      & 0.85 & \texttt{solid}          & Muted; completed \\
\texttt{at-risk}     & Amber         & 1.0  & \texttt{solid}          & Warning; attention needed \\
\texttt{blocked}     & Red           & 1.0  & \texttt{solid}          & Danger; cannot proceed \\
\texttt{cancelled}   & Grey-light    & 0.60 & \texttt{diagonal-hatch} & De-emphasized; struck through \\
\texttt{tentative}   & Grey-light    & 0.70 & \texttt{dashed-border}  & Uncertain; may not happen \\
\bottomrule
\end{tabular}
\caption{Status-to-style map structure (colour roles; each theme supplies hex values)}
\label{tab:status-style}
\end{table}

\paragraph{Pattern vocabulary.}
Patterns are defined as named SVG \texttt{<pattern>} elements in the theme's \texttt{<defs>}
block. All themes must use the same pattern geometry so that cross-theme pattern rendering
is consistent:
\begin{itemize}
  \item \texttt{solid} --- no pattern; fill and stroke colours only.
  \item \texttt{diagonal-hatch} --- 45° lines at 4~px spacing, 0.5~px stroke width,
        stroke colour = \texttt{status.stroke}, fill = \texttt{status.fill} at 30\% opacity.
  \item \texttt{dashed-border} --- solid fill at \texttt{status.fill}; border rendered as a
        dashed stroke (4~px on, 4~px off) at \texttt{status.stroke}.
\end{itemize}

\paragraph{Status-map completeness.}
A theme with a missing status entry is malformed. The renderer must detect and report this
at theme load time, not at render time.

\subsubsection{Category-to-Style Map}
\label{sec:category-style-map}

Categories are arbitrary strings defined by IR authors. Themes provide an optional
\texttt{category\_map} that overrides colour for activities carrying a matching
\texttt{category} value:

\begin{lstlisting}[language=yaml,caption={Category map example}]
category_map:
  infrastructure: { fill: "#2D6A8F", stroke: "#1A4A6A" }
  security:       { fill: "#8B0000", stroke: "#5A0000" }
  ux:             { fill: "#4A7C59", stroke: "#2A5C39" }
\end{lstlisting}

Category styling overrides \emph{only} the \texttt{fill} and \texttt{stroke} of the status
entry. Opacity and pattern are always drawn from the status map. This preserves the semantic
visual signal of status (done $=$ muted, cancelled $=$ hatched) while allowing brand-aligned
colour coding per initiative type.

If an activity's category is absent from \texttt{category\_map}, status-only styling applies
without error or warning.

\subsection{Built-in Themes}
\label{sec:builtin-themes}

Five built-in themes address the primary use cases identified in David's research. Each is a
complete, self-contained style specification; none requires IR changes or additional IR
fields.

\subsubsection{Consulting Theme}
\label{sec:theme-consulting}

\textbf{Visual identity.}
McKinsey/BCG-style executive roadmap. Black, white, and a single accent colour (default
navy \texttt{\#1F497D}). Heavy typography, generous whitespace, no decorative chrome.
Optimised for A4/Letter printing and A3 poster output. Visual language matches consulting
deliverables that non-technical executives expect~\cite{thinkcell}.

\begin{itemize}
  \item \textbf{Font:} Helvetica Neue (Arial fallback); weights 400/600/700. Embedded for
        metric determinism.
  \item \textbf{Canvas:} 1200~px wide; 48~px top/bottom margins; track header width 180~px.
  \item \textbf{Axis:} No gridlines. Tick marks only. Quarter labels in bold uppercase
        (e.g., \textbf{Q1\,2026}). Today marker disabled by default.
  \item \textbf{Tracks:} Wide headers (180~px). Thick separator (1~px, \texttt{\#333333}).
        No row background colour alternation.
  \item \textbf{Activities:} Square corners (\texttt{bar\_radius: 0}). Two status colours:
        navy for planned/in-progress, mid-grey for done. Cancelled activities are rendered
        at 0\% opacity with diagonal hatching at 20\% opacity (structural presence without
        visual weight). No progress bar by default.
  \item \textbf{Milestones:} Solid black diamonds, 14~px half-diagonal. Label right, bold,
        9~pt.
  \item \textbf{Legend:} Suppressed (\texttt{position: none}). Status is communicated by
        colour; a legend would clutter the clean aesthetic.
\end{itemize}

\textbf{Use case:} Programme timelines, transformation roadmaps, board presentations. The
deliberate sparseness focuses attention on the strategic narrative rather than operational
detail.

\subsubsection{Executive Theme}
\label{sec:theme-executive}

\textbf{Visual identity.}
Corporate boardroom aesthetic. Navy/slate palette on white, with subtle alternating quarter
column shading and a gentle border radius on bars. Status is communicated by both colour and
a small status icon at the left of each bar. Designed for 16:9 slide dimensions.

\begin{itemize}
  \item \textbf{Font:} Georgia or Garamond (serif headings); Helvetica Neue (body). Both
        embedded.
  \item \textbf{Canvas:} 1600~px wide (16:9 target); 40~px margins; header width 160~px.
  \item \textbf{Axis:} Light grey gridlines (0.5~px). Month abbreviation labels. Alternating
        quarter shading (\texttt{section\_band\_alpha: 0.04}). Today marker enabled
        (red dashed line, ``Today'' label).
  \item \textbf{Tracks:} Medium headers (150~px). Alternating row background: white /
        \texttt{\#F8F8F8}. Thin separator (0.5~px, \texttt{\#DDDDDD}).
  \item \textbf{Activities:} Rounded corners (\texttt{bar\_radius: 4}). Full status palette:
        navy (planned), royal blue (in-progress), silver (done), amber (at-risk), coral
        (blocked), grey hatched (cancelled), grey dashed-border (tentative). Status icon
        (8~pt glyph) at bar left edge; icon colour is white on dark fills.
  \item \textbf{Milestones:} Navy diamonds with a white centre dot (4~px radius inset fill).
        Label right, 9~pt.
  \item \textbf{Legend:} Right-aligned, full status legend.
\end{itemize}

\textbf{Use case:} Quarterly business reviews, investor roadmaps, product strategy
presentations. Polished without being sparse; the status palette lets executives absorb
delivery health at a glance.

\subsubsection{Product Theme}
\label{sec:theme-product}

\textbf{Visual identity.}
Developer-friendly; information-dense. Per-track colour coding from the \texttt{color} hint
field. Category colours take visual precedence over status colours; status is communicated
by a small icon rather than fill. Progress bars always visible when set.

\begin{itemize}
  \item \textbf{Font:} Inter or Roboto (system-friendly sans-serif), 10~pt base. Compact
        sizing for high information density.
  \item \textbf{Canvas:} 1400~px wide; tighter margins (32~px top/bottom); header 140~px.
  \item \textbf{Axis:} Dashed monthly gridlines. Today marker always visible and prominent.
        Week numbers as secondary axis labels when \texttt{axis\_unit = week}.
  \item \textbf{Tracks:} Color-coded header backgrounds when \texttt{track.color} hint is
        set. Compact rows (36~px). Sub-lanes shown without track-height expansion cue ---
        the density is expected.
  \item \textbf{Activities:} Slightly rounded bars (\texttt{bar\_radius: 3}). Category
        colours override status fill; status glyph icon at bar left edge. Progress bar
        always rendered (even if thin) when \texttt{progress} is set.
  \item \textbf{Milestones:} Smaller diamonds (10~px half-diagonal). Label \emph{below}
        diamond, 8~pt. Stacks horizontally when space permits.
  \item \textbf{Legend:} Bottom; full category $+$ status legend in two rows.
\end{itemize}

\textbf{Use case:} Product roadmaps for engineering and product-management audiences.
Dense information; appropriate for teams who need per-track status at a glance.

\subsubsection{Release Theme}
\label{sec:theme-release}

\textbf{Visual identity.}
Engineering release calendar. Traffic-light status palette (green/amber/red). Monochrome
headers. Today marker is the most prominent element. Optimised for CI/CD dashboard display
and dense sprint/release schedules.

\begin{itemize}
  \item \textbf{Font:} JetBrains Mono or Courier New (monospace). Consistent character
        widths simplify label truncation computations.
  \item \textbf{Canvas:} 1200~px wide; tight margins (24~px top/bottom); header 120~px.
  \item \textbf{Axis:} Solid weekly gridlines (1~px, \texttt{\#CCCCCC}). Today marker:
        bold red solid line (2~px), ``TODAY'' label in uppercase red above axis.
  \item \textbf{Tracks:} No separator lines; alternating row backgrounds (\texttt{\#FFFFFF}
        / \texttt{\#F4F4F4}) imply boundaries. Minimal header (dark text on light bg).
  \item \textbf{Activities:} No border radius (\texttt{bar\_radius: 0}). Traffic-light fill:
        \texttt{\#2E7D32} (done), \texttt{\#1565C0} (in-progress), \texttt{\#757575}
        (planned), \texttt{\#F57F17} (at-risk), \texttt{\#C62828} (blocked); diagonal
        hatching for cancelled; dashed border for tentative. No status icon; colour is
        the sole status signal.
  \item \textbf{Milestones:} Upward-pointing triangle (shape: \texttt{triangle-up}),
        12~px height, filled with status colour. This is a theme-level shape override;
        the IR still uses the Milestone entity type.
  \item \textbf{Legend:} Top-right; compact icon $+$ one-word label per status.
\end{itemize}

\textbf{Use case:} Software release schedules, deployment windows, sprint calendars.
Engineers over executives; clarity of status over aesthetic refinement.

\subsubsection{Minimal Theme}
\label{sec:theme-minimal}

\textbf{Visual identity.}
Maximum whitespace, minimum decoration. All bars are a single dark grey regardless of status;
status is communicated by pattern alone (hatching for cancelled, dashed border for tentative).
Print-optimised; renders correctly in greyscale. Designed to blend into reports and LaTeX
documents without visual dominance.

\begin{itemize}
  \item \textbf{Font:} Times New Roman or Latin Modern (LaTeX-compatible serif). 10~pt base.
  \item \textbf{Canvas:} 900~px wide; narrow margins (24~px top/bottom); header 140~px.
  \item \textbf{Axis:} No gridlines. No tick marks. Period labels only (year or quarter),
        in 9~pt regular weight. No today marker.
  \item \textbf{Tracks:} No header background. Single thin separator (0.5~px,
        \texttt{\#AAAAAA}). No row alternation.
  \item \textbf{Activities:} No border radius. All bars \texttt{\#333333} fill, regardless
        of status. Pattern differentiates cancelled (diagonal-hatch) and tentative
        (dashed-border). No progress bar. No icons. Done activities rendered at 65\% opacity.
  \item \textbf{Milestones:} Open diamonds (stroke only, no fill), 1~px stroke. Label right,
        9~pt regular.
  \item \textbf{Legend:} None (\texttt{position: none}).
\end{itemize}

\textbf{Use case:} Academic papers, internal reports, LaTeX-embedded timelines, any context
where the diagram must be subordinate to surrounding text.

\subsection{Cross-Theme Visual Comparison}
\label{sec:theme-comparison}

To illustrate data-independence, consider a single IR with three tracks (Platform, Mobile,
Data), six activities across statuses planned / in-progress / done / at-risk / cancelled /
tentative, and two milestones (Beta, GA), spanning 2026-Q1 through 2026-Q4.

\begin{table}[h]
\centering
\small
\begin{tabularx}{\textwidth}{lXXXXX}
\toprule
\textbf{Property} & \textbf{Consulting} & \textbf{Executive} & \textbf{Product} & \textbf{Release} & \textbf{Minimal} \\
\midrule
Background       & White        & White            & White           & White           & White \\
Axis labels      & Bold Q labels & Small-caps months & Week numbers   & Week dates      & Year only \\
Gridlines        & None         & Light quarterly  & Dashed monthly  & Solid weekly    & None \\
Track header     & 180~px bold  & 150~px serif     & 140~px colored  & 120~px mono     & 140~px serif \\
Bar corners      & Square       & Rounded 4~px     & Rounded 3~px    & Square          & Square \\
Status signal    & Colour only  & Colour $+$ icon  & Icon $+$ colour & Colour only     & Pattern only \\
Done bar         & Grey         & Silver muted     & Grey + icon     & Green           & Dark 65\% opacity \\
Cancelled bar    & 0\% opacity  & Grey hatch       & Grey hatch      & Grey hatch      & Hatch \\
Progress bar     & Hidden       & Shown            & Always shown    & Hidden          & Hidden \\
Milestones       & Filled black & Navy $+$ dot     & Small, below    & Triangle        & Open diamond \\
Today marker     & Off          & Red dashed       & Red solid       & Bold red        & None \\
Legend           & None         & Right, full      & Bottom, full    & Top-right       & None \\
\bottomrule
\end{tabularx}
\caption{Visual comparison: same IR rendered under five built-in themes}
\label{tab:theme-comparison}
\end{table}

\paragraph{Geometric invariant.}
All five renderings share \emph{identical} geometry: the same bar widths, milestone
$x$-positions, track heights, and label coordinate values. Only style differs. This
invariant can be verified by asserting coordinate equality across the five outputs before
applying theme styling.

\subsection{Theme Inheritance and Extension}
\label{sec:theme-inheritance}

Themes support single-level inheritance via the \texttt{extends} key:

\begin{lstlisting}[language=yaml,caption={Custom theme extending Consulting}]
theme_id:     "acme-corp"
display_name: "ACME Corporate"
extends:      "consulting"

canvas:
  background_color: "#F5F5F0"   # warm paper white

typography:
  font_family: "Futura, Century Gothic, sans-serif"
  font_files:
    regular: "fonts/Futura-Book.woff2"
    bold:    "fonts/Futura-Bold.woff2"

status_map:
  planned:
    fill:   "#005A8E"   # ACME primary brand blue
    stroke: "#003A5E"
  in-progress:
    fill:   "#003A5E"
    stroke: "#001A2E"
\end{lstlisting}

Inheritance rules:
\begin{enumerate}
  \item Properties not specified in the child are taken from the parent exactly.
  \item \texttt{status\_map} and \texttt{category\_map} are merged entry-by-entry: a child
        overrides only the entries it declares; unspecified entries inherit from parent.
  \item Introducing a new font family requires supplying the corresponding embedded font
        files. A theme that names a font without providing files is malformed.
  \item Inheritance depth is exactly one level. Chains (\texttt{A} extends \texttt{B}
        extends \texttt{C}) are not supported; each theme must be resolvable from its
        parent alone.
  \item A theme may \texttt{extend: "default"} explicitly to inherit from the baseline
        default theme without inheriting any of the five named built-in themes.
\end{enumerate}

\subsection{Theme-Engine Contract (Formal)}
\label{sec:theme-contract}

A conforming theme engine must satisfy all of the following:

\begin{enumerate}
  \item \textbf{Accept all valid IR.} A theme engine must not reject a well-formed IR
        document for content-related reasons. An unrecognised \texttt{status} value emits
        a warning and falls back to \texttt{planned} styling.
  \item \textbf{Require no additional IR fields.} The theme engine reads only the fields
        defined in Section~\ref{sec:ir}. It may read hint fields (\texttt{activity.color},
        \texttt{activity.icon}) as advisory inputs but must not fail if they are absent.
  \item \textbf{Full status-map coverage.} The theme's \texttt{status\_map} must contain
        entries for all seven IR status values. A missing entry is a theme authoring error
        detectable at theme-load time.
  \item \textbf{Geometry immutability.} The theme engine applies visual treatment only; it
        may not alter any coordinate computed by the rendering phases in
        Section~\ref{sec:layout-algorithm}. Style does not affect geometry.
  \item \textbf{Hint fields are advisory.} The \texttt{activity.color} and
        \texttt{activity.icon} hint fields may be silently ignored by the theme engine. If
        honoured, they supplement status styling (colour override only); they do not replace
        the pattern or opacity components of the status-map entry.
  \item \textbf{Unknown categories fall back silently.} If an activity's \texttt{category}
        is absent from the theme's \texttt{category\_map}, status-only styling applies
        without error or warning. New categories never require theme changes; they are
        transparent to the theme engine until explicitly mapped.
\end{enumerate}
